% Generated by scripts/materialize_unified_paper.py; do not edit.
\section{Rigidity of exact gradient merge}\label{low:sec:rigidity}

For a target ordered set \([r]\), consider deterministic triangular weights
\(W_{ij}\) and the connection-plus-mean potential
\[
F_i(b)=\sum_{j>i}\log\Phi(-b_j)
-\frac1d\sum_{i<j<q}W_{jq}m(b_j)m(b_q).
\tag{L3.1}\label{low:main:eq:3-1}
\]
At the central cutoff,
\[
 \partial_jF_i(c\mathbf1)
 =-m_0-\frac{m_0m'_0}{d}
 \sum_{\substack{q>i\\q\ne j}}W_{jq}.
\tag{L3.2}\label{low:main:eq:3-2}
\]

\begin{proposition}[Unique exact-merge array]\label{low:prop:rigidity}
Normalize the coefficient of
\(\sum_{j>i}\log\Phi(-b_j)\) in \eqref{low:main:eq:3-1} to be one.  If the linear term in
\eqref{low:main:eq:3-2}, multiplied by
\(b_j-c=\sqrt d\,h_{ij}\), must equal \(-W_{ij}h_{ij}\) at every boundary,
then the weights are uniquely determined by
\[
 W_{ij}=m_0\sqrt d+
 \frac{m_0m'_0}{\sqrt d}
 \sum_{\substack{q>i\\q\ne j}}W_{jq}.
\tag{L3.3}\label{low:main:eq:3-3}
\]
There is no further normalized scalar degree of freedom.
\end{proposition}

\begin{proof}
The coefficient comparison gives \eqref{low:main:eq:3-3}.  Process the first endpoint in
decreasing order.  Every weight on the right side has already been defined,
so the system is triangular and has a unique solution.
\end{proof}

\begin{remark}[Global presentation scale]\label{low:rem:global-scale}
Multiplying the entire inequality by \(\kappa>0\) is a presentation
equivalence, not a new normalized array.  After dividing the connection term
back to coefficient one, the normalized array still satisfies \eqref{low:main:eq:3-3}.  If one
instead calls \(W^{\rm eff}=\kappa W\) the effective coefficient without
renormalizing, then its affine recursion has base term
\(\kappa m_0\sqrt d\), not \(m_0\sqrt d\).  Thus a normalized solution and a
\(\kappa\)-scaled effective coefficient are never identified.  No assertion
is made about edge-dependent rescalings.
\end{remark}

The feedback operator in \eqref{low:main:eq:3-3} has row norm at most
\(\rho=m_0m'_0/D=o(1)\).  Its accumulated increment is
\[
 2A_0\binom r4[1+O(\rho)].
\tag{L3.4}\label{low:main:eq:3-4}
\]
With the source centered coefficient \(m_0^2[1+o(1)]\), the exact-merge
ledger contributes the previously reviewed term
\((1+o(1))/(96\log C)\).  Exact feedback evaluation can change only the
next order.
